\section{Problem Setup}

\subsection{Reed--Solomon parameters}

Let \(\F\) be a finite field and let
\(\mathcal L=(a_i)_{i\in I}\) be an injectively indexed evaluation domain,
where \(|I|=n\).  For \(1\le k\le n\), define
\[
  \mathcal C=\RS[\F,\mathcal L,k]
  =
  \left\{
    (f(a_i))_{i\in I}:
    f\in\F[X],\ \deg f<k
  \right\}.
\]
The rate is \(\rho=k/n\).  The smooth-domain regime takes \(\rho\) from a fixed
constant set, such as
\[
  \left\{\frac12,\frac14,\frac18,\frac1{16}\right\},
\]
and requires \(\mathcal L\) to have an appropriate smooth structure.

\subsection{Relative Hamming distance}

For a finite alphabet \(A\) and \(u,v\in A^I\), let
\[
  d_H(u,v)=|\{i\in I:u_i\ne v_i\}|,
  \qquad
  \Delta(u,v)=\frac{d_H(u,v)}{n}.
\]
For a nonempty code \(\mathcal D\subseteq A^I\), define
\[
  \Delta(w,\mathcal D)
  =
  \min_{c\in\mathcal D}\Delta(w,c).
\]
Unless stated otherwise, radii are relative radii in \([0,1]\).
The uniform MCA bound also allows arbitrary real values of its radius
parameter. Natural-number radii describe the finite Hamming grid and its
rounding; distances in the main statements are normalized by \(n\).

\subsection{Interleaving}

For a positive integer \(m\), define
\[
  \mathcal C^{\equiv m}
  =
  \left\{
    p\in(\F^m)^I:
    \text{for every }j\in\{1,\ldots,m\},
    \ (p_{i,j})_{i\in I}\in\mathcal C
  \right\}.
\]
Hamming distance on \((\F^m)^I\) counts a coordinate \(i\) as erroneous when
the complete vector symbol differs.  It does not sum componentwise scalar
distances.

\subsection{List-decoding balls}

For \(w\in(\F^m)^I\) and \(\delta\in[0,1]\), define
\[
  \List(\mathcal C^{\equiv m},w,\delta)
  =
  \left\{
    c\in\mathcal C^{\equiv m}:
    \Delta(w,c)\le\delta
  \right\}.
\]
The worst-case list size is
\[
  L_m(\mathcal C,\delta)
  =
  \max_{w\in(\F^m)^I}
    |\List(\mathcal C^{\equiv m},w,\delta)|.
\]
The maximum over received words is essential.  A bound on a structured subset
of received words is an application result, not the universal list-decoding
statement.
